% --- IMPORTS ---
\documentclass[leqno]{article}
\usepackage{graphicx}
\usepackage{dsfont}
\usepackage{mathtools}
\usepackage{amssymb}
\usepackage{amsmath}
\usepackage{amsthm}
\usepackage{float}
\usepackage{ragged2e}
\usepackage{tensor}
\usepackage{fancyhdr}
\usepackage{empheq}
\usepackage[most]{tcolorbox}
\usepackage{enumitem}
\usepackage{dirtytalk}
\usepackage{marginnote}
\usepackage{microtype}
\usepackage[
  a4paper,
  left=0.8in,
  right=1in,
  top=1in,
  bottom=0.5in,
  headheight=14pt,
  headsep=0.4in
]{geometry}
\setlength{\parindent}{0cm}
% --- TikZ Packages ---
\usepackage{pgfplots}
\pgfplotsset{compat=1.18}
\usepgfplotslibrary{fillbetween, polar}
\usetikzlibrary{calc, positioning}
\usetikzlibrary{angles, quotes}
\usetikzlibrary{arrows.meta}
\usetikzlibrary{decorations.pathreplacing, decorations.markings}
\usetikzlibrary{patterns}
\usetikzlibrary{intersections}
% --- URL Packages ---
\usepackage[colorlinks=true,linkcolor=red,citecolor=magenta,urlcolor=black]{hyperref}%
\urlstyle{same}
% --- END OF IMPORTS ---

% --- CUSTOM COMMANDS ---
\RenewDocumentCommand{\marks}{o m}{%
  \IfValueTF{#1}%
    {\marginnote{\raggedright\textbf{[#2]}}[#1]}%
    {\marginnote{\raggedright\textbf{[#2]}}}%
}
\newcommand{\vspacerule}[2]{%
  \par\vspace*{#1}%
  \noindent\hrulefill\par
  \vspace*{#2}%
}
\definecolor{scarlet}{RGB}{205, 40, 75}
\newcommand{\questionitem}{%
  \item\phantomsection
  \addcontentsline{toc}{section}{Question \number\value{enumi}}%
}
\renewcommand{\contentsname}{Questions}
% --- END OF CUSTOM COMMANDS ---

% --- HEADER CONFIG ---
\pagestyle{fancy}
\fancyhead{}
\fancyfoot{}
\renewcommand{\headrulewidth}{0.9pt}
\lhead{Matrix Determinants and Inverses}
\chead{}
\rhead{\href{https://danieljsmith.org}{danieljsmith.org}}
% --- END OF HEADER CONFIG ---

\begin{document}
\vspace*{20pt}
\tableofcontents
\newpage
\begin{enumerate}[
  label=\textbf{\arabic*.},
  leftmargin=*,
  topsep=0pt,
  partopsep=0pt,
  parsep=0pt
]

% ---- Question 1 ----
\questionitem

\[
B = \begin{pmatrix} 1 & a \\ 1 & 3 \end{pmatrix}
\] \\
where $a$ is a real constant and $a \neq 3$. \\

\begin{enumerate}[label=\textbf{(\alph*)}]
  \item Find $B^{-1}$ in terms of $a$. \marks{3} \\
\end{enumerate}

Given that
\[
B^{-1} = 2\mathbf{I} - \frac{1}{2}B
\] \\
where $\mathbf{I}$ is the $2 \times 2$ identity matrix, \\

\begin{enumerate}[label=\textbf{(\alph*)}, resume]
  \item find the value of $a$. \marks{3} \\
\end{enumerate}


\newpage

% ---- Question 2 ----
\questionitem

Given that
\[
A=\begin{pmatrix}k+1 & 2\\ 1-2k & k-4\\ 2k-1 & 4-k\end{pmatrix}
\quad \text{and} \quad
B=\begin{pmatrix}1 & 1 & 1\\ 1 & 0 & 1\end{pmatrix}
\] \\
where $k$ is a constant, \\

\begin{enumerate}[label=\textbf{(\alph*)}]
  \item determine the matrix $BA$. \marks{2} \\
  \item hence find the values of $k$ for which $BA$ is singular. \marks{3} \\
\end{enumerate}

\newpage

% ---- Question 3 ----
\questionitem

A community centre offers three evening courses: art, dance and pottery. At the start of the year there were $580$ enrolments in total. The number enrolled in art was $20$ more than the total number enrolled in dance and pottery combined. \\

By the end of the year:
\begin{itemize}
\item the number enrolled in art had decreased by $5\%$
\item the number enrolled in dance had increased by $10\%$
\item the number enrolled in pottery had increased by $15\%$
\item overall enrolment had increased by $20$
\end{itemize}
Form and solve a matrix equation to find the initial number enrolled in each course. \marks{7} \\

\newpage

% ---- Question 4 ----
\questionitem

The matrix $A$ is given by
\[
A=\begin{pmatrix}
2 & k & 1\\
1 & 0 & 2\\
k & -1 & 1
\end{pmatrix}
\] \\
where $k$ is a real constant. \\

\begin{enumerate}[label=\textbf{(\alph*)}]
  \item Show that $A$ is non-singular for all real values of $k$. \marks{3} \\
  \item Given that
  \[
\renewcommand{\arraystretch}{1.5}
A^{-1}=\begin{pmatrix}
  \tfrac{1}{2} & -\frac{1}{2} & \frac{1}{2}\\
  \frac{1}{4} & \frac{1}{4} & -\frac{3}{4}\\
  -\frac{1}{4} & \frac{3}{4} & -\frac{1}{4}
\end{pmatrix}
  \]
  find the value of $k$. \marks{2} \\
\end{enumerate}

\newpage

% ---- Question 5 ----
\questionitem

A cinema chain sells three types of annual pass, $A$, $B$ and $C$. \\

It is known that pass $A$ makes a profit of \pounds 16 per pass sold, pass $B$ makes a profit of \pounds 15 per pass sold and pass $C$ makes a profit of \pounds 12 per pass sold. \\

Between 2023 and 2024
\begin{itemize}
\item sales of pass $A$ increased by $25\%$
\item sales of pass $B$ decreased by $20\%$
\item sales of pass $C$ increased by $50\%$
\end{itemize}

In total $900\,000$ passes were sold in 2023. \\

In total $1\,050\,000$ passes were sold in 2024. \\

The total profit from these passes in 2024 was \pounds 15.12 million. \\

Form and solve a matrix equation to find, to $2$ significant figures, the number of each type of pass sold in 2023. \marks{8} \\

\newpage

% ---- Question 6 ----
\questionitem

\[
A = \begin{pmatrix}
a & 1 & 0 \\
2 & 3 & 1 \\
0 & 1 & a
\end{pmatrix}
\] \\
where $a$ is a real constant. \\

\begin{enumerate}[label=\textbf{(\alph*)}]
  \item Determine the values of $a$ for which $A$ is singular. \marks{2} \\
  \item Hence, for all other values of $a$, find $A^{-1}$ in terms of $a$. \marks{4} \\
\end{enumerate}

\newpage

% ---- Question 7 ----
\questionitem

You are given that $a$ is a real parameter. \\

The matrix $A$ is given by
\[
A=\begin{pmatrix}
1 & -2 & 3 \\
4 & 1 & -1 \\
a & 0 & 1
\end{pmatrix}
\] \\

\begin{enumerate}[label=\textbf{(\alph*)}]
  \item Find an expression for $\det(A)$ in terms of $a$. \marks{2} \\
\end{enumerate}

You are given the following system of equations in $x$, $y$ and $z$. \\

\[
\begin{aligned}
x-2y+3z&=-6 \\
4x+y-z&=7 \\
ax+z&=a-1
\end{aligned}
\] \\

The system can be written in the form
\[
A\begin{pmatrix}x \\ y \\ z\end{pmatrix}=\begin{pmatrix}-6 \\ 7 \\ a-1\end{pmatrix}
\] \\

\begin{enumerate}[label=\textbf{(\alph*)}, resume]
  \item
  \begin{enumerate}[label=(\roman*)]
    \item In the case where $A$ is not singular, solve the given system by using $A^{-1}$. \marks{5}
    \item In the case where $A$ is singular, describe the configuration of the planes whose equations are the three equations of the system. \marks{3} \\
  \end{enumerate}
\end{enumerate}

The transformation represented by $A$ is denoted by $T$. \\
A solid of volume $12$ is transformed by $T$ to an image in 3-D space. \\

\begin{enumerate}[label=\textbf{(\alph*)}, resume]
  \item
  \begin{enumerate}[label=(\roman*)]
    \item Determine the range of values of $a$ for which the orientation of the image is the reverse of the orientation of the object. \marks{1}
    \item Determine the range of values of $a$ for which the volume of the image is less than the volume of the object. \marks{2} \\
  \end{enumerate}
\end{enumerate}

\newpage


\newpage

% ---- Question 8 ----
\questionitem

\[
A = \begin{pmatrix}
k & 2 & 1 \\
-1 & 3 & 0 \\
4 & -2 & 1
\end{pmatrix}
\qquad
B = \begin{pmatrix}
3 & -4 & -3 \\
1 & k - 4 & -1 \\
-10 & 2k + 8 & 3k + 2
\end{pmatrix}
\] \\
where $k$ is a constant. \\

\begin{enumerate}[label=\textbf{(\alph*)}]
  \item Determine the value of the constant $c$ for which
  \[
  AB = (3k - c)I
  \]
  \marks[-20pt]{2}
  \item Hence determine the value of $k$ for which $A^{-1}$ does not exist. \marks{2} \\
\end{enumerate}

Given that $A^{-1}$ does exist, \\

\begin{enumerate}[label=\textbf{(\alph*)}, resume]
  \item write down $A^{-1}$ in terms of $k$. \marks{1} \\
  \item Use the answer to part (c) to solve the simultaneous equations
  \[
  \begin{aligned}
  kx + 2y + z &= 1 \\
  -x + 3y &= 2 \\
  4x - 2y + z &= 3
  \end{aligned}
  \]
  giving the values of $x$, $y$ and $z$ in simplest form in terms of $k$. \marks{3} \\
\end{enumerate}

\newpage

% ---- Question 9 ----
\questionitem

A company manufactures only three models of bicycle: Commuter, Electric and Folding. \\

Each bicycle manufactured is one of these three models. \\

In 2023, a total of $1000$ bicycles were manufactured. \\

The number of Folding bicycles manufactured was $200$ fewer than the combined number of Commuter and Electric bicycles. \\

In 2024 the number manufactured of
\begin{itemize}
\item Commuter increased by $10\%$
\item Electric decreased by $5\%$
\item Folding increased by $5\%$
\end{itemize}
In 2024, the total number of bicycles manufactured increased by $3.8\%$. \\

\begin{enumerate}[label=\textbf{(\alph*)}]
  \item
  \begin{enumerate}[label=(\roman*)]
    \item Define, for each model, a variable for the number of bicycles manufactured in 2023. \marks{2}
    \item Using your variables from part (a)(i), write down three equations that model this situation. \marks{2}
  \end{enumerate}
  \item By forming and solving a matrix equation, determine how many bicycles of each model were manufactured in 2023. \marks{4} \\
\end{enumerate}

\newpage

% ---- Question 10 ----
\questionitem

Three planes have equations
\[
\begin{aligned}
x+2y-z&=4,\\
3x+ky+2z&=k+11,\\
ky+4z&=k+4
\end{aligned}
\] \\
where $k$ is a constant. \\

\begin{enumerate}[label=\textbf{(\alph*)}]
  \item Show that the planes have a unique point of intersection except for one value of $k$, which should be found. \marks{4} \\
  \item Show that, whenever the planes have a unique point of intersection, this point is independent of $k$. Find the point. \marks{6} \\
\end{enumerate}
\end{enumerate}
\end{document}
